% !TeX spellcheck = it_IT
\newpage
\section{Introduzione}

\subsection{Problemi, Algoritmi e Programmi}

\begin{definition}[Problema]
    Un \textbf{problema} descrive un compito generale che deve essere risolto. Non si riferisce a un particolare input, ma a una classe di possibili input uniti agli output considerati corretti.
\end{definition}

\begin{definition}[Problema computazionale]
    Un \textbf{problema computazionale} è un problema le cui istanze e output possono essere rappresentati in una forma che può essere manipolata da un modello computazionale, come un programma o una macchina di Turing. Più specificatamente, un problema computazionale definisce:
    \begin{itemize}
        \item un insieme di istanze ammissibili;
        \item 
    \end{itemize}
\end{definition}



\subsection{Lessico, Sintassi e Semantica}

\subsection{Architettura del compilatore}

\subsection{Principio di Composizionalità}

\subsection{Linguaggi formali e Grammatiche}
